% Compile with: uplatex file.tex && dvipdfmx file.dvi
%           or: latexmk -uplatex file.tex
%\documentclass[uplatex,dvipdfmx]{jsarticle}
\documentclass[dvipdfmx,a4paper,11pt]{jsarticle}
\usepackage{amsmath, amssymb}
\usepackage{tikz}

\title{scenario\_transverse\_field\_ising}
\date{\today}

\begin{document}
\maketitle

2サイト横磁場イジングモデルを扱う。結合定数 $J$（イジング相互作用の強さ）と横磁場の強さ $\Gamma$ をスカラーパラメータとして準備する。

J, $\Gamma$ をスカラーパラメータとして定義して

\[
  J, \Gamma
\]

両サイトで使うスピン演算子 $S$ を準備する。

S をスピン演算子として定義して

\[
  S, S^{x}, S^{y}, S^{z}
\]

横磁場イジングハミルトニアン $H=-JS^z_1S^z_2-\Gamma(S^x_1+S^x_2)$（イジング結合項と横磁場項）を登録する。

$H = -J Sz_1 Sz_2 - \Gamma (Sx_1 + Sx_2)$

\begin{equation}
  H = -J Sz_1 Sz_2 - \Gamma (Sx_1 + Sx_2)
  \label{eq:E1}
\end{equation}

ハミルトニアン $H$ とサイト1の $S^z$ の交換子を計算し、SU(2)交換関係に基づく運動方程式の一部を確認する。

[E1, $Sz_1$]

\begin{equation}
  2 \Gamma \mathit{i} Sy_1
  \label{eq:E2}
\end{equation}

平均場近似を行うため、サイト2のスピンを置き換える平均場（c-数）$m_z=\langle S^z\rangle$ をスカラーパラメータとして準備する。

mz をスカラーパラメータとして定義して

\[
  mz
\]

サイト2の演算子 $S^z_2$ だけを平均場 $m_z$ に置き換え、平均場ハミルトニアン $H_{\rm MF}$ を得る（サイト1の演算子や横磁場項は変化しないはずである）。

E1 の $S^z_2$ を mz に置き換えて

\begin{equation}
  H = - \Gamma Sx_1 - \Gamma Sx_2 - J mz Sz_1
  \label{eq:E3}
\end{equation}

\end{document}